\section{A direct global-to-local obstruction for the full smooth-cutoff equation}
\label{sec:direct-projection-intertwiner-obstruction}

\paragraph{Review status.}
The main theorem, printed-twist corollary and invariant-domain
form have complete post-writing author and independent
readings, each with its exact separately recorded scope.

The preceding local-to-global results impose only the exact
inside/outside plateaus. A different, purely operator-theoretic
argument applies in the global-to-local direction when the
\emph{full} physical-cutoff equation is required. It needs no
inverse and no information about zeros. The full equation is
an explicit hypothesis of this theorem, not an equation
attributed to the unprinted comparison section of Masullo.

For the rational local source, with $S=\{\infty\}$ and the
printed twist $A=1/2$, retain the exact completion
\[
 H_s^{\rm loc}=L^2(\mathbb R,p_s(x)\,dx),\qquad
 p_s(x)=(1+x^2)^{s/2},\qquad s\ge0.
 \tag{MA-PI.230}
\]
Its inverse additive realization is
\[
 (V^{-1}g)(t)=\frac12|t|^{-1/2}g(\log|t|)
                  \quad(t\ne0),
 \qquad
 \|g\|_{H_s^{\rm loc}}^2
       =2\int_{\mathbb R}|V^{-1}g(t)|^2
                         p_s(\log|t|)\,dt.
 \tag{MA-PI.231}
\]
The unspecified value at $t=0$ is irrelevant in $L^2$.
In particular, $V^{-1}:H_s^{\rm loc}\to L^2(\mathbb R)^{\rm even}$
is bounded, with norm at most $1/\sqrt2$.
The normalized actions satisfy
\[
 V^{-1}\tau_y=D_yV^{-1},\qquad
 (D_yf)(t)=e^{-y/2}f(e^{-y}t),\qquad
 \widehat{D_yf}(\xi)=e^{y/2}\widehat f(e^y\xi).
 \tag{MA-PI.232}
\]
All these formulas hold first pointwise off zero and then
in their actual Hilbert spaces. They use the mass-one
rational quotient Haar normalization established in MA-PI-010.

Fix just one positive physical level $Y$ and one of the
original smooth real profiles, so that
\[
 \begin{gathered}
  \rho\in C^\infty([0,\infty)),\qquad 0\le\rho\le1,\\
  \rho(r)=1\quad(0\le r\le Y/2),\qquad
  \rho(r)=0\quad(r\ge2Y),\\
  B=\mathcal F^{-1}M_{\rho(|\xi|)}\mathcal F
           \quad\hbox{on }L^2(\mathbb R)^{\rm even}.
 \end{gathered}
 \tag{MA-PI.233}
\]
This is the literal physical multiplier of MA-PI.5 in the
rational $S=\{\infty\}$ sector. Its $L^2$ norm is at most
one; no boundedness assertion in the weighted local norm
is made or required.

\paragraph{MA-PI-011 (no nonzero bounded exact reverse intertwiner).}
Let $\mathscr J$ be a Hilbert space of global logarithmic
vectors, invariant under every translation $\tau_y$, with
each translation bounded. Let $P:\mathscr J\to\mathscr J$
be a bounded idempotent. In particular, $\mathscr J$ may be
the actual closed centered theta ideal and $P$ the global
orthogonal PW projection onto a closed subspace of that ideal.
Suppose a bounded linear map
\[
 U:\mathscr J\longrightarrow H_s^{\rm loc},\qquad
 U\tau_y=\tau_y U\quad(y\in\mathbb R)
 \tag{MA-PI.234}
\]
satisfies the full equation
\[
       B V^{-1}U h=V^{-1}U P h
 \tag{MA-PI.235}
\]
for $h$ in a dense linear subspace of $\mathscr J$.
Then $U=0$. The conclusion holds for every $s\ge0$.
Neither an inverse for $U$, a smooth output, an
invariant testing core, nor any zero-location premise is needed.

\begin{proof}
Put $T=V^{-1}U$. By (MA-PI.231), $T$ is bounded from
$\mathscr J$ to ordinary even $L^2(\mathbb R)$. Both sides
of (MA-PI.235) are therefore continuous in $h$ with values
in that ordinary $L^2$ space. The dense-core equation extends
to every $h\in\mathscr J$, without an assumption that the
core is invariant under $P$.

Since $P^2=P$ and $P$ maps $\mathscr J$ to itself, the
extended equation can be used twice:
\[
       B^2Th=BTP h=TP^2h=TP h=BTh.
 \tag{MA-PI.236}
\]
Every operation here is legitimate in ordinary $L^2$.
The full equation supplies weighted local membership of
the actual outputs when it is needed; it does not assert
that $B$ is bounded on all of $H_s^{\rm loc}$.
Taking the ordinary $L^2$ Fourier transform gives
\[
       q(\xi)\widehat{Th}(\xi)=0\quad\hbox{a.e.},\qquad
       q(\xi)=\rho(|\xi|)^2-\rho(|\xi|).
 \tag{MA-PI.237}
\]

The smooth profile has both a one plateau and a zero
plateau. By continuity there is a nonempty open interval
$(a,b)\subset(0,\infty)$ on which $0<\rho<1$.
Thus $q$ is nonzero on $(a,b)$ and on $(-b,-a)$.
For each $y\in\mathbb R$, apply (MA-PI.237) to
$\tau_yh\in\mathscr J$. Equivariance and (MA-PI.232)
give
\[
       q(\xi)e^{y/2}\widehat{Th}(e^y\xi)=0
                                      \quad\hbox{a.e.}
 \tag{MA-PI.238}
\]
Consequently $\widehat{Th}$ vanishes almost everywhere
on both $e^y(a,b)$ and $e^y(-b,-a)$.
It suffices to take the countable set $y\in\mathbb Q$:
these intervals cover $\mathbb R\setminus\{0\}$, since
for each positive $r$ the interval
$(\log(r/b),\log(r/a))$ contains a rational number.
The countable union of the exceptional null sets is still
null. Hence $\widehat{Th}$ is an $L^2$ function supported,
up to a null set, at the singleton $\{0\}$. It is zero.
The $L^2$ Fourier isometry yields $Th=0$, and injectivity
of $V^{-1}$ yields $Uh=0$. This holds for every $h$.
\end{proof}

\paragraph{The full printed range of twists.}
The same conclusion holds in every printed local norm
$0<A<1$, after the exact half-centering is retained.
Put $b=A-1/2$, so that $-1/2<b<1/2$, and replace
$H_s^{\rm loc}$ by
\[
 \begin{gathered}
 H_{b,s}^{\rm loc}=L^2(\mathbb R,e^{2bx}p_s(x)\,dx),\\
 \|Vf\|_{H_{b,s}^{\rm loc}}^2
       =2\int_{\mathbb R}|f(t)|^2|t|^{2b}
                              p_s(\log|t|)\,dt.
 \end{gathered}
 \tag{MA-PI.239}
\]
The chart $Vf=2e^{x/2}f(e^x)$ and the normalized actions
remain exactly (MA-PI.231)--(MA-PI.232). In the literal
printed unweighted $p_s$ realization, multiply the local
logarithmic output by $e^{bx}$; the corresponding action
is $e^{by}\tau_y$, as in MA-PI-007. This change is the
exact unitary, not an arbitrary equivalence of norms.

Suppose $U:\mathscr J\to H_{b,s}^{\rm loc}$ is bounded,
linear and translation-equivariant, and (MA-PI.235) holds
on a dense core, now as an equality of tempered
distributions in the additive coordinate. Then again
$U=0$, with no zero-location input.

\begin{proof}
Let $w(t)=|t|^{2b}p_s(\log|t|)$. For $f=V^{-1}g$ and
$\varphi\in\mathcal S(\mathbb R)$, Cauchy--Schwarz gives
\[
 |\langle f,\varphi\rangle|
 \le\frac{\|g\|_{H_{b,s}^{\rm loc}}}{\sqrt2}
     \left(\int_{\mathbb R}
          |\varphi(t)|^2|t|^{-2b}
                         p_s(\log|t|)^{-1}\,dt\right)^{1/2}.
 \tag{MA-PI.240}
\]
The last integral is bounded by a fixed Schwartz
seminorm squared. Near zero, $p_s^{-1}\le1$ and
$2b<1$ make $|t|^{-2b}$ integrable. At infinity, a
sufficient polynomial Schwartz bound is integrable.
Thus $V^{-1}:H_{b,s}^{\rm loc}\to\mathcal S'(\mathbb R)$
is a continuous injective map, represented by a locally
integrable even function. This estimate also justifies
the additive dilation identities in distributions.

Fourier transform and multiplication by the fixed smooth
compact function $\rho(|\xi|)$ are continuous on
$\mathcal S'$. The dense-core equation therefore extends
to every $h\in\mathscr J$ in $\mathcal S'$. Idempotence
gives (MA-PI.237) as a distributional identity. Applying
it to all translated inputs gives (MA-PI.238) in the
same sense. On every open transition interval the
multiplier $q$ is nonzero, so division by $q$ on compact
subintervals proves that $\widehat{Th}$ vanishes there.
Dilations of those intervals cover all nonzero
frequencies. Consequently
\[
                 \operatorname{supp}\widehat{Th}\subset\{0\}.
 \tag{MA-PI.241}
\]

A tempered distribution supported at one point has
finite order there and is a finite sum of derivatives
of the delta distribution. For completeness, continuity
gives one finite Schwartz-seminorm bound; the Taylor
cutoff argument of MA-PI.202--203, in its scalar form,
then annihilates every test with the corresponding
finite origin jets zero. Prescribing those jets gives
the stated finite delta expansion. Inverse Fourier
transform therefore makes $Th$ a polynomial.
If this polynomial is nonzero of degree $n$, its
weighted squared norm at positive infinity dominates
a positive multiple of
\[
             \int_R^\infty t^{2n+2b}\,dt=\infty,
                    \qquad n\ge0,\quad b>-1/2.
 \tag{MA-PI.242}
\]
Here $p_s\ge1$; thus even degree zero is excluded.
This contradicts (MA-PI.239). Hence $Th=0$, and $U=0$.
\end{proof}

\paragraph{Invariant-domain form (MA-PI-011).}
The boundedness hypothesis has one precise role above:
it extends a dense-core equation. If the equation is
already imposed on an invariant domain, the algebraic
obstruction requires no continuity assumption.
Let $\mathscr D$ be a vector space with an action
$S_y$ of $(\mathbb R,+)$, let $P:\mathscr D\to\mathscr D$
be an idempotent, and let
\[
 T:\mathscr D\longrightarrow\mathcal S'(\mathbb R)^{\rm even}
 \quad\hbox{be linear},\qquad
 T S_y=D_yT,\qquad BT=TP.
 \tag{MA-PI.243}
\]
Here $B$ is the one fixed smooth physical multiplier
in (MA-PI.233), and $D_y$ is the normalized additive
dilation in (MA-PI.232), acting also on distributions.
Then every $Th$ is a polynomial. If
\[
             T(\mathscr D)\cap\mathbb C[t]=\{0\},
 \tag{MA-PI.244}
\]
it follows that $T=0$. In particular, the polynomial-free
codomain may be the actual even flat Schwartz local
source, or any of the additive Hilbert realizations in
(MA-PI.239). No boundedness or continuity of $T$ or $P$
is required in this invariant-domain version.

\begin{proof}
Every displayed operation is defined on the stated
domain because $P(\mathscr D)\subset\mathscr D$.
The two applications of $BT=TP$ give
$B^2Th=TP^2h=TPh=BTh$. Fourier transformation yields
$q\widehat{Th}=0$. Apply the same identity to $S_yh$
for every $y$ and use the exact dilation covariance.
Division by the nonzero smooth multiplier on the
translated transition intervals proves
$\operatorname{supp}\widehat{Th}\subset\{0\}$,
exactly as in (MA-PI.241). The scalar finite-order
point-support argument already proved above makes $Th$
a polynomial. Assumption (MA-PI.244) gives zero.
Even Schwartz space contains no nonzero polynomial;
the weighted Hilbert cases were checked in (MA-PI.242).
\end{proof}

This statement does not assert that the global orthogonal
PW projector preserves the original smooth theta ideal.
For a proposed smooth comparison, its actual domain
must be stable under both $P$ and every $S_y$, and the
full equation must be defined on that domain. Neither
of these hypotheses follows from density or from an
identity of scalar traces.

\paragraph{Exact source interface and scope.}
The local smooth space $E$ on Masullo p.~515 is not itself
a Hilbert space. The completion used here is the separately
specified local line norm of Definition~10.1.1.1(2),
(10.1.1.4)--(10.1.1.5), p.~434, at $A=1/2$ and $I=\{0\}$.
The corollary uses the same exact printed norm at every
$0<A<1$, still with $I=\{0\}$.
The finite sign quotient in Definition~2.3.1.1,
(2.3.1.2), pp.~99--100, and the Haar definitions on
pp.~556,560--562 give exactly (MA-PI.231).
The physical frequency ball is the first bullet of
Step~II on p.~516, with its inverse-Fourier composition
in (12.4.1.8), p.~517. The global orthogonal PW source
is Definition~11.2.1.3, (11.2.1.7)--(11.2.1.8),
pp.~488--489. Its range is a closed subspace of the actual
global theta ideal, so its restriction to that ideal
is an idempotent endomorphism. The source interfaces and
their reading scopes are recorded separately from this proof.

This theorem directly excludes nonzero bounded maps in
the stated global-to-local class, including possible
isomorphisms. It uses only one fixed smooth transition
profile and equivariance under all normalized dilations.
It does not address sharp characteristic-function cutoffs,
mere plateau implications, or identities holding only
after a limit, a Gram correction, or a change of action.
No nonzero transition interval exists for a sharp cutoff,
and the idempotence step is unavailable for mere plateaus.
These distinctions prevent an assertion that every possible
prime-bearing comparison is impossible.

The theorem is independent of MA-PI-008's critical-line
input. It supplies no arithmetic vanishing of the FT002
deficiency and no new statement of RH. Its categorical
content is that an exact equivariant morphism must respect
the idempotent operation on actual completed objects;
the smooth physical multiplier cannot do so nontrivially
in this class. The inherited comparison-map and nuclear-pairing
placement remains cote~1, scans~079--080 and~111--112,
before trace cyclicity or Fredholm determinant operations.
Finite Artin local systems and their ordinary equivariant
maps retain the SGA~1, Expos\'e~V, \S7, and Deligne (1979),
Definition~6.1 and \S\S6.2--6.3 locators. No positivity,
purity or norm assertion follows from a $K_0$ or
$\lambda$-ring class alone.
